\section{问题重述}

\subsection{问题背景}

大语言模型的训练同时消耗数据与算力。参数规模和训练数据量共同影响验证集交叉熵损失（Loss），模型在不同任务上的能力还须通过多维基准评测衡量。已有标度律揭示了规模与损失的经验联系；在算力受限时，仅依据参数量和数据量难以确定资源应如何分配。

训练语料的质量与领域构成使配置问题更为复杂。质量指标的方向和评价角度各异，同一文本可能得到相互矛盾的评价；不同领域的训练占比也会改变损失表现。质量提升和长上下文处理占用额外算力，与基础训练共享预算。即使 Loss 下降，其对应的下游能力变化仍需另行衡量。因此，有限预算下的资源配置需要同时考虑数据因素、规模效应、计算成本与能力评测。

\subsection{问题提出}

\noindent\textbf{问题一：}训练语料的质量评价涉及 22 项方向各异的指标，同一文本还可能出现显著不一致的指标判断。利用附件 A 抽样集及两个扩展集的全部记录，需建立样本质量评分及其领域聚合，比较抽样集与扩展集的域级评价，并核查评分与文本内容的一致性；对指标冲突应说明成因和消解后的综合评价依据，在扩展集考察其适用性。附件 A 的配比实验还要求刻画 17 个训练领域的配比及其组合与 Loss 的关系，其中各配比非负且总和为 1。该关系须经指定检验集验证，并结合外推表讨论稳健性；如纳入质量评分，还应说明两套领域划分的对应依据。

\noindent\textbf{问题二：}经典标度律未显式反映质量和领域配比。依据问题一的输出及附件 B 的训练轨迹，建立以参数规模、训练数据量、质量和配比解释 Loss 的广义关系，并检验其在不同模型族、文献数据及较大参数范围内的适用性。附件 A、B 的实验相互独立，损失口径的可比性须说明；质量相关的半合成资料和大模型估算值也应与直接观测区分。模型应给出参数估计和验证结果，揭示各因素的边际效用与弹性、领域数据间的替代或互补关系，并说明质量提升能够替代参数扩张的条件。

\noindent\textbf{问题三：}在前两问关系的支持下，以总算力预算为约束，确定参数规模、训练数据量、质量及领域配比的配置，使预测 Loss 尽可能低。成本包括基础训练、质量提升和长文本注意力开销，三者默认使用相同的训练数据量；领域配比可参与寻优，也可依据前两问结果给定，但须解释选择。至少比较三个不同数量级的预算，并考察质量成本函数变化对配置的影响。对预算变化时可能发生的结构性转移，应给出数学定义与识别依据。上下文长度由附件 C 的架构数据确定可行范围，作为外生条件；须解析求得注意力与基础训练开销相当的临界长度，并分析可行长度变化对配置的影响。

\noindent\textbf{问题四：}利用附件 C 的多维评测、时间序列及模型计算量、训练数据量等资料，确定综合能力的度量方式，并对逐任务评测结果作聚合分析。须建立动力学或因果推断模型，分别量化规模扩张与非规模技术进步对能力增长的贡献占比；研究对象的开放权重与许可证等筛选口径、预训练与对话或微调模型的区分、时间轴口径均应明确。在算力增长放缓的情景下，预测未来 12 个月或 24 个月开源模型的能力前沿并分析不确定性。前三问的 Loss 与本问的多维评测得分还需借助分级可比的桥接数据或可核验文献建立联系，评价映射误差对结论的影响。
